%======================================================================
% CAPÍTULO 01: INICIO DE SESIÓN
%======================================================================
% Proceso de autenticación y primer acceso
%======================================================================

\chapter{Inicio de Sesión}
\label{ch:inicio-sesion}

\section{Presentación del Sistema}

El \textbf{SSSIFANB} (Sistema de Seguridad Social de la Fuerza Armada Nacional Bolivariana) es una plataforma tecnológica integral diseñada para gestionar los procesos de afiliación, administración de datos y otorgamiento de prestaciones sociales dirigidas al personal militar y sus beneficiarios.

El sistema fue desarrollado por el \textbf{Instituto de Previsión Social de la Fuerza Armada Nacional Bolivariana (IPSFANB)} con el objetivo de garantizar la transparencia, eficiencia y trazabilidad en la administración de los recursos destinados a la protección social del personal militar.

\subsection*{Marco Legal}

El SSSIFANB opera bajo el marco normativo establecido por:

\begin{itemize}[leftmargin=*]
    \item \textbf{Constitución de la República Bolivariana de Venezuela} (Art. 328): Determina que la Seguridad Social es la base de la suprema felicidad pública.
    \item \textbf{Ley del Sistema de Seguridad Social de la FANB} (Decreto Gaceta Oficial N° 6.209): Ley Negro Primero, promulgada para dignificar al profesional militar.
    \item \textbf{Leyes y Reglamentos Complementarios}
\end{itemize}

\begin{figure}[H]
\centering
\includegraphics[width=0.9\textwidth]{1-presentacion.png}
\caption{Pantalla de presentación del SSSIFANB}
\label{fig:presentacion}
\end{figure}

\section{Pantalla de Login}

El acceso al sistema se realiza a través de la pantalla de autenticación:

\begin{figure}[H]
\centering
\begin{minipage}{0.48\textwidth}
\centering
\includegraphics[width=\textwidth]{1.1-login.png}
\caption{Pantalla de inicio de sesión}
\label{fig:login}
\end{minipage}
\hfill
\begin{minipage}{0.48\textwidth}
\centering
\includegraphics[width=\textwidth]{1.2-validacion-2FA.png}
\caption{Validación 2FA}
\label{fig:2fa}
\end{minipage}
\end{figure}

\subsection*{Campos del Formulario}

\begin{longtable}{p{3cm}p{9cm}}
\toprule
\textbf{Campo} & \textbf{Descripción} \\
\midrule
\endhead
\bottomrule
\caption{Campos del formulario de login}
\label{tab:campos-login}
\endlastfoot

\textbf{Usuario} & Cédula de identidad del afiliado \\
\textbf{Clave} & Contraseña personal de acceso \\
\textbf{Código 2FA} & Código de 6 dígitos (si está habilitado) \\
\end{longtable}

\section{Autenticación de Dos Factores}

Para mayor seguridad, el sistema implementa autenticación de dos factores (2FA/TOTP):

\begin{enumerate}[leftmargin=*]
    \item Acceda a la configuración desde su perfil
    \item Seleccione la opción de habilitar 2FA
    \item Escanee el código QR con su aplicación autenticadora
    \item Ingrese el código de verificación
    \item Guarde los códigos de respaldo
\end{enumerate}

\begin{tcolorbox}[
    colback=sandraSoft,
    colframe=sandraGreen,
    width=0.9\textwidth,
    halign=center
]
\small
\textbf{Tip de Seguridad}: Nunca comparta sus códigos de verificación. El IPSFANB nunca solicitará su contraseña por teléfono o correo.
\end{tcolorbox}

\section{Pantalla Principal}

Una vez autenticado, el usuario accede a la pantalla principal del sistema:

\begin{figure}[H]
\centering
\begin{minipage}{0.48\textwidth}
\centering
\includegraphics[width=\textwidth]{2-principal.png}
\caption{Pantalla principal}
\label{fig:pantalla-principal}
\end{minipage}
\hfill
\begin{minipage}{0.48\textwidth}
\centering
\includegraphics[width=\textwidth]{2.1-principal-configuracion.png}
\caption{Acceso a configuración}
\label{fig:config-acceso}
\end{minipage}
\end{figure}

La pantalla principal incluye:
\begin{itemize}[leftmargin=*]
    \item \textbf{Barra superior}: Logo, nombre del usuario, notificaciones
    \item \textbf{Menú lateral}: Navegación a los módulos (Afiliación, Beneficiarios, TIM, Prestaciones, Nómina)
    \item \textbf{Área de contenido}: Información del módulo seleccionado
    \item \textbf{Pie de página}: Información institucional
\end{itemize}

\section{Configuración de Usuario}

\begin{figure}[H]
\centering
\begin{minipage}{0.48\textwidth}
\centering
\includegraphics[width=\textwidth]{2.2-configuracion.png}
\caption{Panel de configuración}
\label{fig:config-panel}
\end{minipage}
\hfill
\begin{minipage}{0.48\textwidth}
\centering
\includegraphics[width=\textwidth]{2.3-cambiar-clave.png}
\caption{Cambio de contraseña}
\label{fig:cambiar-clave}
\end{minipage}
\end{figure}

\subsection*{Cambio de Contraseña}

Requisitos de contraseña:
\begin{itemize}[leftmargin=*]
    \item \textbf{Longitud mínima}: 8 caracteres
    \item \textbf{Mayúsculas y minúsculas}: Al menos una de cada
    \item \textbf{Números}: Al menos un dígito
    \item \textbf{Caracteres especiales}: Al menos un símbolo
\end{itemize}

\begin{figure}[H]
\centering
\begin{minipage}{0.48\textwidth}
\centering
\includegraphics[width=\textwidth]{2.4-permitir-2FA.png}
\caption{Habilitar 2FA}
\label{fig:permitir-2fa}
\end{minipage}
\end{figure}

\section{Cierre de Sesión}

Para cerrar sesión de manera segura:

\begin{enumerate}[leftmargin=*]
    \item Haga clic en su nombre de usuario
    \item Seleccione \textbf{Cerrar Sesión}
    \item Confirme la acción
\end{enumerate}

El sistema invalidará el token y lo redirigirá a login.

\section{Arquitectura Técnica del Sistema}

El SSSIFANB está desarrollado con las siguientes tecnologías:

\begin{longtable}{p{3cm}p{9cm}}
\toprule
\textbf{Componente} & \textbf{Tecnología} \\
\midrule
\endhead
\bottomrule
\caption{Tecnologías del sistema}
\label{tab:tecnologias}
\endlastfoot

\textbf{Frontend} & Angular CLI v17+ (Standalone Components) \\
\textbf{Runtime} & Node.js 18.x \\
\textbf{Capa de datos} & API RESTful hacia Sandra Server \\
\textbf{Autenticación} & OAuth2 con Tokens JWT \\
\end{longtable}

\subsection*{Estructura de Directorios}

\begin{verbatim}
src/
├── app/
│   ├── core/           # Servicios core (API, interceptors, interfaces)
│   ├── feature/        # Módulos de interacción (pages, layouts)
│   └── shared/         # Componentes reutilizables
\end{verbatim}

\newpage
